\KOMAoption{paper}{a4, portrait}
\areaset{267mm}{180mm}
\fancyheadoffset{0pt} % обновляем хедер и футер

\newgeometry{
  top=17mm,
  bottom=10mm,
  right=8mm,
  left=20mm,
}

\begin{landscape}

\setlength{\headsep}{7.5mm}
\setlength{\topmargin}{-27mm}
\setlength{\footskip}{21pt}

\color{uniblue}{\section[(обязательное) Подключение дискретных входов \newline и выходов устройства к внутренним цепям алгоритмов для архитектуры II типа]{(обязательное)\\Подключение дискретных входов и выходов устройства к внутренним цепям алгоритмов для архитектуры II типа}\label{app:inouts2}}
\color{black}

\begin{figure}[h!]
  \centering
  
\includegraphics[width=1.1\textwidth]{\apppath/Приложение. Подкл вх и вых2/_latex/img/1.pdf}
  \caption{Подключение дискретных входов и выходов к внутренним сигналам алгоритмов для плат в слоте М1}
  \label{fig:sig21} % Опционально: ссылка на рисунок
\end{figure}

\begin{figure}[h!]
  \centering
  \includegraphics[width=1.5\textwidth]{\apppath/Приложение. Подкл вх и вых2/_latex/img/2.pdf}
  \caption{Подключение дискретных входов и выходов к внутренним сигналам алгоритмов для плат в слотах М3,М4}
  \label{fig:sig22} % Опционально: ссылка на рисунок
\end{figure}

\begin{figure}[h!]
  \centering
  \includegraphics[width=1.45\textwidth]{\apppath/Приложение. Подкл вх и вых2/_latex/img/3.pdf}
  \caption{Подключение дискретных входов и выходов к внутренним сигналам алгоритмов для плат в слотах М5-М7}
  \label{fig:sig23} % Опционально: ссылка на рисунок
\end{figure}

\begin{figure}[h!]
  \centering
  \includegraphics[width=1.5\textwidth]{\apppath/Приложение. Подкл вх и вых2/_latex/img/4.pdf}
  \caption{Подключение дискретных входов и выходов к внутренним сигналам алгоритмов для плат в слоте М8}
  \label{fig:sig24} % Опционально: ссылка на рисунок
\end{figure}

\clearpage
\end{landscape}
\restoregeometry


